% Original vector schematic of the axial-resolution geometry (not to scale).
% Two on-axis emitters, one at the image plane and one displaced axially, seen from
% the central lenslet and from the outermost retained lenslet.
\begin{tikzpicture}[
    x=1cm,y=1cm,
    font=\footnotesize,
    ink/.style={text=black!85},
    lens/.style={<->,>={Stealth[length=2mm]},draw=black!75,line width=0.8pt},
    rayA/.style={draw=blue!60!black,line width=0.8pt},
    rayB/.style={draw=red!65!black,line width=0.8pt,dash pattern=on 3pt off 2pt},
    plane/.style={draw=black!45,dashed,line width=0.5pt},
    dim/.style={<->,>={Stealth[length=1.6mm]},draw=black!70,line width=0.5pt},
    lab/.style={font=\footnotesize,ink,inner sep=1.5pt,align=center}
]
% ---- geometry (cm; not to scale) ----
% image plane at x=0, Fourier lens at x=8, emitter E1 displaced by 1.6 behind the image plane.
\def\xf{8}      % Fourier lens position (f_FL)
\def\dz{1.6}    % axial displacement at the image plane, exaggerated
\def\hs{3.0}    % ray height at the FL for the outermost view (d_max)
\pgfmathsetmacro{\tn}{\hs/\xf}           % tan(theta')
\pgfmathsetmacro{\hB}{(\xf+\dz)*\tn}     % ray height at the FL for E1
\pgfmathsetmacro{\sh}{\dz*\tn}           % lateral shift at the image plane
\pgfmathsetmacro{\ang}{atan(\tn)}

% optical axis = line of sight of the central view
\draw[draw=green!45!black,line width=0.9pt] (-\dz-0.6,0) -- (\xf,0);

% image plane plane and Fourier lens
\draw[plane] (0,-0.3) -- (0,2.4);
\draw[lens] (\xf,-4.0) -- (\xf,4.3);

% sub-aperture of the outermost retained lenslet on the FL
\draw[draw=orange!80!black,line width=2.6pt] (\xf,2.5) -- (\xf,3.9);

% lines of sight of the outermost view (parallel, angle theta')
\draw[rayA] (0,0) -- (\xf,\hs);
\draw[rayB] (-\dz,0) -- (\xf,\hB);

% emitters
\fill[blue!60!black] (0,0) circle (0.07);
\fill[red!65!black] (-\dz,0) circle (0.07);
\node[lab,anchor=north west,text=blue!60!black] at (0.1,-0.12) {focused\\emitter};
\node[lab,anchor=north,text=red!65!black] at (-\dz,-0.12) {defocused\\emitter};

% lateral shift s seen at the image plane in the outer view
\fill[red!65!black] (0,\sh) circle (0.06);
\draw[draw=black!55,line width=0.4pt] (0,\sh) -- (-2.55,1.55);
\node[lab,anchor=east] at (-2.5,1.6) {$s=\Delta z_{im}\tan\theta'$};

% angle theta'
\draw[draw=black!70,line width=0.5pt] (1.9,0) arc (0:\ang:1.9);
\node[lab,anchor=west] at (1.95,0.28) {$\theta'$};

% dimensions
\draw[dim] (-\dz,-1.0) -- (0,-1.0);
\node[lab,anchor=north] at (-\dz/2,-1.05) {$\Delta z_{im}=M^{2}\,\Delta z$};
\draw[dim] (0,-1.9) -- (\xf,-1.9);
\node[lab,fill=white] at (4,-1.9) {$f_{FL}$};
\draw[dim] (\xf+0.75,0) -- (\xf+0.75,\hs);
\node[lab,anchor=west] at (\xf+0.82,\hs/2) {$d_{max}$};
\node[lab,anchor=south,text=orange!80!black] at (\xf-0.25,4.4)
    {outermost retained\\lenslet (sub-aperture)};

% names
\node[lab,anchor=north] at (\xf,-4.05) {Fourier lens};
\node[lab,anchor=north,text=green!45!black] at (5.0,-0.18)
    {central lenslet: both emitters at the same point};

% sub-image profiles in the outer view (lateral axis magnified)
\def\xi{12.6}
\draw[->,>={Stealth[length=1.6mm]},draw=black!70,line width=0.5pt]
    (\xi-0.3,-1.0) -- (\xi-0.3,3.6);
\node[lab,anchor=east] at (\xi-0.35,3.4) {lateral};
\draw[draw=blue!60!black,line width=0.8pt,samples=60,domain=-0.7:0.7]
    plot ({\xi+1.9*exp(-2.7726*(\x)^2/0.36)},{2.6*\x+0.4});
\draw[draw=red!65!black,line width=0.8pt,dash pattern=on 3pt off 2pt,
    samples=60,domain=-0.1:1.3]
    plot ({\xi+1.9*exp(-2.7726*(\x-0.6)^2/0.36)},{2.6*\x+0.4});
\draw[dim] (\xi+2.35,0.4) -- (\xi+2.35,1.96);
\node[lab,anchor=west] at (\xi+2.42,1.18) {$s$};
\draw[dim] (\xi+0.95,-0.38) -- (\xi+0.95,1.18);
\draw[draw=black!55,line width=0.4pt] (\xi+0.95,-0.38) -- (\xi+1.15,-0.7);
\node[lab,anchor=west] at (\xi+1.15,-0.7) {$R_{im}$ (FWHM)};
\node[lab,anchor=south] at (\xi+0.9,4.0)
    {sub-image in the\\outer view (at the limit)};
\end{tikzpicture}
